\chapter*{Conclusions of Part~III}
\addcontentsline{toc}{chapter}{Conclusions of Part~III}
\label{ch:part3-concl}

The three chapters in this Part develop the neutrino-induced phase shift into a comprehensive observational program spanning CMB power spectra and 21-cm cosmology.

In Chapter~\ref{ch:neutrino-phase-shift}, we introduced two complementary template-based methods for isolating the phase shift in CMB data. The spectrum-based template measures the coherent displacement of acoustic peaks across multipoles, while the perturbation-based template captures the shift directly at the level of photon-baryon perturbations where it is originally imprinted. Applying both methods to current CMB datasets from \textit{Planck}, ACT, and SPT, we detected the neutrino-induced phase shift at up to $14\sigma$ significance and confirmed consistency with the Standard Model expectation of three free-streaming neutrino species. The perturbation-based method proved particularly robust across different data releases, and forecasts for the Simons Observatory and CMB-S4 indicate that precision on the phase shift will improve by roughly 50\% and 70\% respectively relative to current measurements.

Chapter~\ref{ch:neutrino-interactions} exploited this framework to place direct constraints on neutrino interactions. By modeling two physically motivated interaction scenarios with scattering rates scaling as $T_\nu^3$ and $T_\nu^5$, we showed that the phase-shift amplitude maps cleanly onto the neutrino decoupling redshift and established lower bounds of $z_{\nu,\mathrm{dec}} > 1.3 \times 10^4$ and $z_{\nu,\mathrm{dec}} > 1.7 \times 10^4$ respectively from combined CMB data. These constraints firmly require that neutrinos were free-streaming well before matter-radiation equality. A notable result is the discovery that even fully fluid-like neutrinos produce a nonzero phase shift at roughly 30\% of the free-streaming value, because even in the tightly coupled limit their sound speed exceeds that of the photon-baryon plasma due to baryon loading. This finding has implications for the interpretation of any future anomalous phase-shift measurement.

In Chapter~\ref{ch:neutrino-21cm}, we extended the phase-shift analysis to the 21-cm brightness temperature during cosmic dawn and identified a qualitatively new feature. The 21-cm signal contains both baryon acoustic oscillations and velocity-induced acoustic oscillations, and the neutrino-induced phase shift in each component differs in both amplitude and scale dependence. The VAO phase shift asymptotes to a value roughly 1.5 times larger than the BAO phase shift and exhibits a smoother transition between large and small scales, reflecting its origin in four-point rather than two-point correlation functions. The combined signature evolves across cosmic dawn as the relative contributions of VAOs and BAOs shift with astrophysical conditions, producing a distinctive redshift-dependent signal that is robust at small scales but sensitive to star formation physics on large scales.

Together, these three chapters establish the neutrino-induced phase shift as a versatile and physically transparent observable that can be tracked across multiple cosmological probes. The progression from a first high-significance detection in the CMB, to direct constraints on neutrino interactions, to the identification of a novel signature in 21-cm cosmology illustrates the broad reach of this approach and its potential to serve as a unified framework for testing neutrino properties with current and upcoming surveys.
